
Or on a par
définition :
@@M0@@
Par ailleurs, la suite exacte :
@@M1@@
donne la relation :
@@M2@@
d'où on déduit :
@@M3@@

Compte tenu de 4.4, la relation à démontrer s'écrit donc :
@@M4@@

Posant @@M5@@, cette relation n'est autre que le théorème de Riemann--Roch appliqué à l'immersion régulière de codimension @@M6@@~@@M7@@, et à l'élément @@M8@@ de @@M9@@.

De plus, comme on a supposé que @@M10@@ vérifie l'hypothèse de (VII 4.1), @@M11@@ est divisible par @@M12@@. On peut donc trouver @@M13@@ tel que @@M14@@ soit encore @@M15@@ (exp. VII 2.7). On est donc ramené au cas où @@M16@@ est un diviseur, et à un élément @@M17@@ tel qu'il existe @@M18@@, avec @@M19@@.

\textbf{4.6.} Il faut donc montrer la relation :
@@M20@@
Or la formule de projection donne :
@@M21@@
et de la suite exacte : @@M22@@, où @@M23@@ est ici inversible, on tire :
@@M24@@
d'où
